\documentclass[10pt,landscape]{article}
\usepackage[utf8]{inputenc}
\usepackage[T1]{fontenc}
\usepackage[brazilian]{babel}
\usepackage[a4paper,landscape,margin=0.55cm]{geometry}
\usepackage[most]{tcolorbox}
\usepackage{xcolor}
\usepackage{tikz}
\usepackage[sfdefault]{cabin}
\usepackage[type1]{comfortaa}
\usepackage{microtype}
\usepackage{ragged2e}
\usepackage{array}
\renewcommand{\familydefault}{\sfdefault}
\pagestyle{empty}
\setlength{\parindent}{0pt}
\setlength{\emergencystretch}{3em}

\definecolor{pmcPurple}{RGB}{176, 92, 214}
\definecolor{pmcOrange}{RGB}{184, 120, 24}
\definecolor{pmcPink}{RGB}{226, 76, 191}
\definecolor{pmcGreen}{RGB}{0, 204, 102}
\definecolor{pmcBlue}{RGB}{58, 179, 216}
\definecolor{textGray}{RGB}{70, 70, 70}

\newcommand{\pmcheader}[2]{%
    \begin{center}
        \begin{tikzpicture}[baseline=(title.base)]
            \node[
                fill=#1,
                text=white,
                rounded corners=0.8mm,
                inner xsep=8pt,
                inner ysep=4pt
            ] (title) {{\comfortaa\fontsize{16}{16}\selectfont\bfseries\itshape #2}};
            \draw[#1, line width=1.1pt] (title.south) -- ++(0, -0.38cm);
        \end{tikzpicture}
    \end{center}
}

\newcommand{\pmcheadercompact}[2]{%
    \begin{center}
        \begin{tikzpicture}[baseline=(title.base)]
            \node[
                fill=#1,
                text=white,
                rounded corners=0.8mm,
                inner xsep=5pt,
                inner ysep=4pt
            ] (title) {{\comfortaa\fontsize{12}{12}\selectfont\bfseries\itshape #2}};
            \draw[#1, line width=1.1pt] (title.south) -- ++(0, -0.38cm);
        \end{tikzpicture}
    \end{center}
}

\newtcolorbox{pmcframe}[1]{
    enhanced,
    colback=white,
    colframe=#1,
    boxrule=0.9pt,
    borderline={1pt}{0pt}{#1, dashed},
    sharp corners,
    arc=0pt,
    left=2.2mm,
    right=2.2mm,
    top=2.2mm,
    bottom=2.2mm
}

\newtcolorbox{pmccard}[3]{
    enhanced,
    colback=white,
    colframe=black!70,
    boxrule=0.8pt,
    sharp corners,
    arc=0pt,
    left=2mm,
    right=2mm,
    top=1.6mm,
    bottom=1.6mm,
    height=#3,
    valign=center,
    before upper={%
        {\raggedright\textbf{\small #1}\par}
        \vspace{0.35em}
        \color{#2}\small\justifying
    }
}

\begin{document}

\begin{center}
    {{\comfortaa\large \bfseries PROJECT MODEL CANVAS}}\\[-0.05cm]
    {\small \textcolor{textGray}{Reconhecimento de Gestos da M\~ao por sEMG com 1 a 2 Sensores no Antebra\c{c}o}}
\end{center}

\vspace{0.1cm}

\noindent
\begin{minipage}[t]{0.19\textwidth}
    \pmcheader{pmcPurple}{POR QU\^E?}
    \vspace{-0.1cm}
    \begin{pmcframe}{pmcPurple}
        \begin{pmccard}{JUSTIFICATIVA}{pmcPurple}{4.05cm}
            H\'a demanda acad\^emica e tecnol\'ogica por interfaces homem-m\'aquina de baixo custo. O prot\'otipo atual com um m\'odulo EMG apresentou alta varia\c{c}\~ao de sinal, exigindo uma solu\c{c}\~ao mais est\'avel e reprodut\'ivel.
        \end{pmccard}
        \vspace{0.18cm}
        \begin{pmccard}{OBJETIVO SMART}{pmcPurple}{3.55cm}
            Projetar e validar, at\'e a etapa final da disciplina, um prot\'otipo com 1 a 2 sensores sEMG capaz de captar gestos simples da m\~ao com maior estabilidade no antebra\c{c}o.
        \end{pmccard}
        \vspace{0.18cm}
        \begin{pmccard}{BENEF\'ICIOS}{pmcPurple}{5.15cm}
            Gera base experimental para classifica\c{c}\~ao futura, reduz ambiguidades de capta\c{c}\~ao, fortalece a aplica\c{c}\~ao embarcada em hardware simples e amplia a reprodutibilidade acad\^emica do sistema.
        \end{pmccard}
    \end{pmcframe}
\end{minipage}
\hfill
\begin{minipage}[t]{0.19\textwidth}
    \pmcheader{pmcOrange}{O QU\^E?}
    \vspace{-0.1cm}
    \begin{pmcframe}{pmcOrange}
        \begin{pmccard}{RESULTADO}{pmcOrange}{3.95cm}
            Um sistema experimental de reconhecimento de gestos da m\~ao por sEMG, com captura no antebra\c{c}o, fixa\c{c}\~ao mec\^anica melhorada e base para comparar 1 sensor versus 2 sensores.
        \end{pmccard}
        \vspace{0.18cm}
        \begin{pmccard}{REQUISITOS}{pmcOrange}{8.95cm}
            Capta\c{c}\~ao est\'avel, posicionamento consistente dos eletrodos, suporte f\'isico em PLA, baixa sensibilidade a ru\'ido e microdeslocamentos, compara\c{c}\~ao entre canais, extra\c{c}\~ao de features simples e viabilidade de processamento embarcado.
        \end{pmccard}
    \end{pmcframe}
\end{minipage}
\hfill
\begin{minipage}[t]{0.39\textwidth}
    \begin{minipage}[t]{0.485\linewidth}
        \pmcheader{pmcPink}{QUEM?}
    \end{minipage}
    \hfill
    \begin{minipage}[t]{0.485\linewidth}
        \pmcheader{pmcGreen}{COMO?}
    \end{minipage}

    \vspace{-0.1cm}

    \begin{minipage}[t]{0.485\linewidth}
        \begin{pmcframe}{pmcPink}
            \begin{pmccard}{STAKEHOLDERS EXTERNOS}{pmcPink}{4.55cm}
                Disciplina e institui\c{c}\~ao de ensino, docente avaliador, laborat\'orio, comunidade acad\^emica e usu\'arios potenciais de interfaces musculares acess\'iveis.
            \end{pmccard}
            \vspace{0.18cm}
            \begin{pmccard}{EQUIPE}{pmcPink}{5.55cm}
                Giuseph D. L. Giangareli, respons\'avel pela revis\~ao bibliogr\'afica, pesquisa de anterioridade, ajustes mec\^anicos da case, montagem, capta\c{c}\~ao e avalia\c{c}\~ao experimental.
            \end{pmccard}
        \end{pmcframe}
    \end{minipage}
    \hfill
    \begin{minipage}[t]{0.485\linewidth}
        \begin{pmcframe}{pmcGreen}
            \begin{pmccard}{PREMISSAS}{pmcGreen}{4.55cm}
                Dois canais em regi\~oes antagonistas tendem a melhorar a separa\c{c}\~ao dos padr\~oes musculares. A case em PLA reduz parte da instabilidade mec\^anica e features no tempo podem atender ao prot\'otipo.
            \end{pmccard}
            \vspace{0.18cm}
            \begin{pmccard}{ESCOPO}{pmcGreen}{5.55cm}
                Revis\~ao do estado da arte, pesquisa de anterioridade, estabiliza\c{c}\~ao f\'isica do m\'odulo, defini\c{c}\~ao do posicionamento dos sensores, coleta inicial e compara\c{c}\~ao entre configura\c{c}\~oes de um e dois sensores.
            \end{pmccard}
        \end{pmcframe}
    \end{minipage}

    \vspace{0.18cm}

    \begin{pmcframe}{pmcGreen}
        \begin{pmccard}{RESTRI\c{C}\~OES}{pmcGreen}{2.2cm}
            Tempo de disciplina, or\c{c}amento reduzido, alta variabilidade anat\^omica, ru\'ido no sinal e disponibilidade limitada de sensores e materiais de fixa\c{c}\~ao.
        \end{pmccard}
    \end{pmcframe}
\end{minipage}
\hfill
\begin{minipage}[t]{0.19\textwidth}
    \pmcheadercompact{pmcBlue}{QUANDO?/QUANTO?}
    \vspace{-0.1cm}
    \begin{pmcframe}{pmcBlue}
        \begin{pmccard}{RISCOS}{pmcBlue}{3.95cm}
            Posicionamento inconsistente dos eletrodos, crosstalk, ru\'ido, pouca repetibilidade entre testes, base de dados pequena e ganho limitado mesmo ap\'os migrar para dois sensores.
        \end{pmccard}
        \vspace{0.18cm}
        \begin{pmccard}{LINHA DO TEMPO}{pmcBlue}{4.55cm}
            Revis\~ao bibliogr\'afica e anterioridade, testes com um sensor, melhoria mec\^anica com case PLA, defini\c{c}\~ao do segundo ponto de capta\c{c}\~ao, compara\c{c}\~ao experimental e consolida\c{c}\~ao final do trabalho.
        \end{pmccard}
        \vspace{0.18cm}
        \begin{pmccard}{CUSTO}{pmcBlue}{3.95cm}
            Proposta de baixo custo, com parte da estrutura reaproveitada. O gasto principal concentra-se em sensores sEMG, eletrodos, fixa\c{c}\~ao, impress\~ao 3D e ajustes para aquisi\c{c}\~ao confi\'avel.
        \end{pmccard}
    \end{pmcframe}
\end{minipage}

\vfill
\begin{center}
    \tiny Fonte: SILVA, Fabiana Big\~ao. \textit{Gerenciamento de projetos fora da caixa}. Rio de Janeiro: Alta Books, 2016.
\end{center}

\end{document}
